\documentclass[a4paper, 11pt]{article}
\usepackage{comment} % enables the use of multi-line comments (\ifx \fi) 
\usepackage{lipsum} %This package just generates Lorem Ipsum filler text. 
\usepackage{fullpage} % changes the margin
\usepackage{libertine}
\usepackage{url}
\linespread{1.1}
\begin{document}
%Header-Make sure you update this information!!!!
\noindent
\large\textbf{Thesis Proposal} \hfill \textbf{Renato Neves} \\
Date: \today

\section*{Simulation of Hybrid Systems via Exact Real Arithmetic}

\subsection*{Motivation and Goals}

Hybrid systems are those that integrate continuous, physical processes, such as
temperature, time, and velocity, with discrete, event-based behaviours, such as
assignments, conditionals, and state
switching~\cite{Platzer10,neves18,goncharov2020}. Examples of such systems
abound, and can be found in diverse critical domains. These
include autonomous driving, modelling of biological systems, and medical and
industrial engineering. It is therefore crucial that one can simulate their
execution as effectively and precisely as possible.

The simulation of hybrid systems, however, still harbours difficult challenges.
A most prominent one arises from the ubiquituous use of \emph{floating-point
arithmetic}~\cite{muller18} -- it introduces errors and potentially leads to
simulations that are completely different from the actual behaviour of the
hybrid system at hand. Consequently the engineer may draw up wrong conclusions
(\emph{e.g.} the driving system is safe) which in turn can cause catastrophic
failures.

The overarching goal of this project is to simulate hybrid systems via
\emph{exact real arithmetic}
instead~\cite{di93,di04,geuvers07,clarke12,bridger19}. More specifically the
idea is to implement a simulator of hybrid systems based on pre-existing
libraries for exact real arithmetic (see for example~\cite{konecny22}). We
expect that the performance of such a simulator will be worse than the standard
counterpart (\emph{i.e.} via floating-point arithmetic), however in principle
it will have much greater fidelity and this is of crucial importance in the
critical domains mentioned above. 

The language for writing down hybrid systems will be based on the one 
in~\cite{goncharov2020}, in a nutshell a while-language~\cite{winskel93} with differential
constructs for specifying how a physical process behaves. The justification for
this choice is that the language is simple, yet powerful, and moreover has a
well-established semantics which we can use as an implementation guide.
The implementation will be based on the programming language \textsc{Haskell},
due to the latter's laziness and higher-order features which is particularly
useful in the context of exact real arithmetic.

\smallskip
\noindent

\subsection*{Research Plan}

The project is divided in three main parts, namely background, implementation, and
benchmarking.

\noindent
\textbf{Background (3 months).}
The project will begin with a brief study on  the fundamentals of exact real
arithmetic~\cite{plume98,geuvers07,bridger19} and with the compilation of a
survey on existing \textsc{Haskell} libraries that support exact real
arithmetic (a good starting point is available in~\cite{konecny22}).

\noindent
\textbf{Implementation (3 months).} 
With the proper bases established, the project then proceeds to the
actual implementation of the simulator of hybrid systems.  As mentioned before, the
implementation will be based on the hybrid language and semantics presented
in~\cite{goncharov2020}.

\noindent
\textbf{Benchmarking (2 months).} 
Finally the project focusses on benchmarking the prospective tool via a
number of carefully chosen case-studies. 

Throughout the project we will register all of our results, conclusions, and
progress in a scientific notebook which will then be compiled into a final
report.

%% Bibliography
\bibliographystyle{alpha}
\bibliography{biblio}

\end{document}
